% GENERATED FILE. Edit canonical lessons, then run npm run generate:books.
% canonical-source-hash: fnv1a64:801631e9960f556b
% canonical-lessons: GU-C389-nirnay, GU-C389-marji, GU-C389-dvidha, GU-C389-khackavu, GU-C389-dagvu

\chapter{Making a Choice}
\label{ch:gu-making-a-choice}
\hlchaptermodality{389}

\begin{chapteropening}
\textbf{By the end of this chapter:} \emph{I can weigh a choice before I decide.}

Complete the last lesson of chapter 389: I can weigh a choice before I decide.
\end{chapteropening}

\section[\texorpdfstring{\gu{નિર્ણય} (nirṇay)}{નિર્ણય (nirṇay)}]{\gu{નિર્ણય} (nirṇay) — a decision}

\label{lesson:GU-C389-nirnay}



\begin{quote}
Before the new one: say the Gujarati for to commence, then the Gujarati for to undertake.
\end{quote}

\subsection*{You'll want to know: \gu{નિર્ણય}}
\textbf{\gu{નિર્ણય}} — \emph{nirṇay} — \textquotedblleft{}a decision\textquotedblright{}.

\begin{quote}
\gu{જો} \gu{તમે} \gu{ઉતાવળમાં} \gu{નિર્ણય} \gu{લેશો} \gu{તો} \gu{પછી} \gu{પસ્તાશો}.
\end{quote}

\begin{quote}
\emph{jo tame utāvaḷmā̃ nirṇay leśo to pachī pastāśo.}
\end{quote}

\begin{quote}
If you decide in a hurry, you will regret it later.
\end{quote}

A masculine noun.

\begin{cousinweb}
\textbf{\gu{નિર્ણય}} (\emph{nirṇay}) is a tatsama, borrowed from Sanskrit \emph{nirṇaya}, \textquotedblleft{}taking off, settling, a decision\textquotedblright{} (Monier-Williams).
\end{cousinweb}

\begin{grammarlens}[title={What you've built}]
One more word for saying what happens if.
\end{grammarlens}

\subsection*{Your turn}
\begin{itemize}
\raggedright
  \item \emph{Say it:} \emph{nirṇay}
  \item \emph{Say it:} \emph{nirṇay}, once more
  \item \emph{Say it:} \emph{jo tame utāvaḷmā̃ nirṇay leśo to pachī pastāśo.} — If you decide in a hurry, you will regret it later.
  \item \emph{Recall:} say the Gujarati for to commence, then the Gujarati for to undertake, then say \emph{nirṇay} again
\end{itemize}

\begin{tcolorbox}[breakable,colback=teal!4,colframe=teal!35!black,title={Before you move on}]
What does \gu{નિર્ણય} mean? (\textquotedblleft{}A decision\textquotedblright{}.) Say it once more.
\end{tcolorbox}
\section[\texorpdfstring{\gu{મરજી} (marjī)}{મરજી (marjī)}]{\gu{મરજી} (marjī) — one's own will}

\label{lesson:GU-C389-marji}



\begin{quote}
Before the new one: say the Gujarati for to undertake, then the Gujarati for a decision.
\end{quote}

\subsection*{You'll want to know: \gu{મરજી}}
\textbf{\gu{મરજી}} — \emph{marjī} — \textquotedblleft{}one's own will\textquotedblright{}.

\begin{quote}
\gu{જો} \gu{તમારી} \gu{મરજી} \gu{હશે} \gu{તો} \gu{આપણે} \gu{કાલે} \gu{જઈશું}.
\end{quote}

\begin{quote}
\emph{jo tamārī marjī haśe to āpaṇe kāle jaīśũ.}
\end{quote}

\begin{quote}
If you are willing, we will go tomorrow.
\end{quote}

A feminine noun.

\begin{cousinweb}
\textbf{\gu{મરજી}} (\emph{marjī}) is usually traced to Arabic \emph{marḍī}, \textquotedblleft{}what pleases\textquotedblright{}, through Persian \emph{marzī}.
\end{cousinweb}

\begin{grammarlens}[title={What you've built}]
One more word for saying what happens if.
\end{grammarlens}

\subsection*{Your turn}
\begin{itemize}
\raggedright
  \item \emph{Say it:} \emph{marjī}
  \item \emph{Say it:} \emph{marjī}, once more
  \item \emph{Say it:} \emph{jo tamārī marjī haśe to āpaṇe kāle jaīśũ.} — If you are willing, we will go tomorrow.
  \item \emph{Recall:} say the Gujarati for to undertake, then the Gujarati for a decision, then say \emph{marjī} again
\end{itemize}

\begin{tcolorbox}[breakable,colback=teal!4,colframe=teal!35!black,title={Before you move on}]
What does \gu{મરજી} mean? (\textquotedblleft{}One's own will\textquotedblright{}.) Say it once more.
\end{tcolorbox}
\section[\texorpdfstring{\gu{દ્વિધા} (dvidhā)}{દ્વિધા (dvidhā)}]{\gu{દ્વિધા} (dvidhā) — a dilemma}

\label{lesson:GU-C389-dvidha}



\begin{quote}
Before the new one: say the Gujarati for a decision, then the Gujarati for one's own will.
\end{quote}

\subsection*{You'll want to know: \gu{દ્વિધા}}
\textbf{\gu{દ્વિધા}} — \emph{dvidhā} — \textquotedblleft{}a dilemma\textquotedblright{}.

\begin{quote}
\gu{જો} \gu{બંને} \gu{નોકરી} \gu{સારી} \gu{હશે} \gu{તો} \gu{મને} \gu{દ્વિધા} \gu{થશે}.
\end{quote}

\begin{quote}
\emph{jo banne nokrī sārī haśe to mane dvidhā thaśe.}
\end{quote}

\begin{quote}
If both jobs are good, I will be in two minds.
\end{quote}

A feminine noun.

\begin{cousinweb}
\textbf{\gu{દ્વિધા}} (\emph{dvidhā}) is a tatsama, borrowed from Sanskrit \emph{dvidhā}, \textquotedblleft{}in two ways, divided\textquotedblright{} (Monier-Williams): a mind split two ways.
\end{cousinweb}

\begin{grammarlens}[title={What you've built}]
One more word for saying what happens if.
\end{grammarlens}

\subsection*{Your turn}
\begin{itemize}
\raggedright
  \item \emph{Say it:} \emph{dvidhā}
  \item \emph{Say it:} \emph{dvidhā}, once more
  \item \emph{Say it:} \emph{jo banne nokrī sārī haśe to mane dvidhā thaśe.} — If both jobs are good, I will be in two minds.
  \item \emph{Recall:} say the Gujarati for a decision, then the Gujarati for one's own will, then say \emph{dvidhā} again
\end{itemize}

\begin{tcolorbox}[breakable,colback=teal!4,colframe=teal!35!black,title={Before you move on}]
What does \gu{દ્વિધા} mean? (\textquotedblleft{}A dilemma\textquotedblright{}.) Say it once more.
\end{tcolorbox}
\section[\texorpdfstring{\gu{ખચકાવું} (khackāvũ)}{ખચકાવું (khackāvũ)}]{\gu{ખચકાવું} (khackāvũ) — to hesitate}

\label{lesson:GU-C389-khackavu}



\begin{quote}
Before the new one: say the Gujarati for one's own will, then the Gujarati for a dilemma.
\end{quote}

\subsection*{You'll want to know: \gu{ખચકાવું}}
\textbf{\gu{ખચકાવું}} — \emph{khackāvũ} — \textquotedblleft{}to hesitate\textquotedblright{}.

\begin{quote}
\gu{જો} \gu{તમારી} \gu{પાસે} \gu{કોઈ} \gu{પ્રશ્ન} \gu{હશે} \gu{તો} \gu{પૂછવામાં} \gu{ન} \gu{ખચકાઓ}.
\end{quote}

\begin{quote}
\emph{jo tamārī pāse koī praśna haśe to pūchvāmā̃ na khackāo.}
\end{quote}

\begin{quote}
If you have a question, do not hesitate to ask.
\end{quote}

Past (he …): \gu{તે} \gu{થોડો} \gu{ખચકાયો}, \emph{te thoḍo khackāyo}.

\begin{cousinweb}
\textbf{\gu{ખચકાવું}} (\emph{khackāvũ}) CDIAL lists the Gujarati verb, \textquotedblleft{}to shrink back\textquotedblright{}, under a reconstructed \emph{khacc}, \textquotedblleft{}sink down or back\textquotedblright{} (CDIAL 3764).
\end{cousinweb}

\begin{grammarlens}[title={What you've built}]
One more word for saying what happens if.
\end{grammarlens}

\subsection*{Your turn}
\begin{itemize}
\raggedright
  \item \emph{Say it:} \emph{khackāvũ}
  \item \emph{Say it:} \emph{khackāvũ}, once more
  \item \emph{Say it:} \emph{jo tamārī pāse koī praśna haśe to pūchvāmā̃ na khackāo.} — If you have a question, do not hesitate to ask.
  \item \emph{Recall:} say the Gujarati for one's own will, then the Gujarati for a dilemma, then say \emph{khackāvũ} again
\end{itemize}

\begin{tcolorbox}[breakable,colback=teal!4,colframe=teal!35!black,title={Before you move on}]
What does \gu{ખચકાવું} mean? (\textquotedblleft{}To hesitate\textquotedblright{}.) Say it once more.
\end{tcolorbox}
\section[\texorpdfstring{\gu{ડગવું} (ḍagvũ)}{ડગવું (ḍagvũ)}]{\gu{ડગવું} (ḍagvũ) — to waver}

\label{lesson:GU-C389-dagvu}



\begin{quote}
Before the new one: say the Gujarati for a dilemma, then the Gujarati for to hesitate.
\end{quote}

\subsection*{You'll want to know: \gu{ડગવું}}
\textbf{\gu{ડગવું}} — \emph{ḍagvũ} — \textquotedblleft{}to waver\textquotedblright{}.

\begin{quote}
\gu{જો} \gu{તમે} \gu{તમારા} \gu{નિર્ણયથી} \gu{ડગશો} \gu{તો} \gu{બીજા} \gu{તમારો} \gu{વિશ્વાસ} \gu{નહીં} \gu{કરે}.
\end{quote}

\begin{quote}
\emph{jo tame tamārā nirṇaythī ḍagśo to bījā tamāro viśvās nahī̃ kare.}
\end{quote}

\begin{quote}
If you waver from your decision, others will not trust you.
\end{quote}

Past (he …): \gu{તે} \gu{એક} \gu{પળ} \gu{માટે} \gu{ડગ્યો}, \emph{te ek paḷ māṭe ḍagyo}.

\begin{cousinweb}
\textbf{\gu{ડગવું}} (\emph{ḍagvũ}) CDIAL lists the Gujarati verb, \textquotedblleft{}to tremble\textquotedblright{}, under a reconstructed \emph{ḍag} (CDIAL 5522); \textquotedblleft{}to waver\textquotedblright{}, a shaking of the mind, is a sense CDIAL does not give.
\end{cousinweb}

\begin{grammarlens}[title={What you've built}]
One more word for saying what happens if.
\end{grammarlens}

\subsection*{Your turn}
\begin{itemize}
\raggedright
  \item \emph{Say it:} \emph{ḍagvũ}
  \item \emph{Say it:} \emph{ḍagvũ}, once more
  \item \emph{Say it:} \emph{jo tame tamārā nirṇaythī ḍagśo to bījā tamāro viśvās nahī̃ kare.} — If you waver from your decision, others will not trust you.
  \item \emph{Recall:} say the Gujarati for a dilemma, then the Gujarati for to hesitate, then say \emph{ḍagvũ} again
\end{itemize}

\begin{tcolorbox}[breakable,colback=teal!4,colframe=teal!35!black,title={Before you move on}]
What does \gu{ડગવું} mean? (\textquotedblleft{}To waver\textquotedblright{}.) Say it once more.
\end{tcolorbox}
